\documentclass{article}
\usepackage{amsthm,amssymb,amsmath}

\newtheorem{problem}{Problem}

\begin{document}
\begin{problem}
Show that any amount of money of at least $\$14$ can be made up
using $\$3$ and $\$8$ coins.
\end{problem}

\begin{proof}

Let $P(M)$: Amount $M$ can be made using $\$3$, and $\$8$ coins,
where $M \ge 14$.

\textbf{Base Case:} For $M=14$, we have

\[
14=1(8)+2(3)
\]

Thus, $P(14)$ holds.

\textbf{Inductive Case:}

Suppose $P(M)$ holds. Then, $M=8a+3b$ for some nonnegative integers
$a$, and $b$.

We shall show that $P(M+1)$ holds by considering two cases.

\textbf{Case $a \ge 1$:}

There is at least one $\$8$ coin. So, $M=8+8(a-1)+3b$. Hence,

\begin{align*}
M+1&=8+8(a-1)+3b+1\\
&=9+8(a-1)+3b\\
&=8(a-1)+3(b+3)
\end{align*}

Since $a-1 \ge 0$ and $b+3 \ge 0$, this expresses $M+1$ as a sum of
$\$3$ and $\$8$ coins.

\textbf{Case $a = 0$:}

Then, $M=3b \ge 14$. Since $b$ is an integer, we have $b \ge 5$. Consequently,
there exists at least five $\$3$ coins. Thus, $M=5(3)+3(b-5)$. Hence,

\begin{align*}
M+1&=5(3)+3(b-5)+1\\
&=16+3(b-5)\\
&=2(8)+3(b-5)
\end{align*}

Since $b-5 \ge 0$, this expresses $M+1$ as a sum of
$\$3$ and $\$8$ coins.
\end{proof}

Therefore, $P(M+1)$ holds for every integer $M \ge 14$. By induction,
every integer $M \ge 14$ can be made using $\$3$ and $\$8$ coins.
\end{document}
